% file: 5.几何向量.tex
% 附录A 几何向量与二次曲面。线性代数教材无此章，内容以同济《高等数学》下册（第八版）
% 第八章「向量代数与空间解析几何」为主要依据，另补入解析几何中教材未展开的少数标准结论
% （点到平面/直线的距离公式、平面束），并写在明确的小节里。
% 只写知识点：定义、定理（写清条件）、核心公式、方法步骤、易错点；不写例题与习题。

\section{向量及其线性运算}

\subsection{向量的概念}

客观世界中有一类量，它们既有大小又有方向，例如位移、速度、加速度、力、力矩，这一类量叫做\textbf{向量}（或矢量）。数学上用一条\textbf{有向线段}表示向量：有向线段的长度表示向量的大小，有向线段的方向表示向量的方向。以 $A$ 为起点、$B$ 为终点的向量记作 $\overrightarrow{AB}$，也常用黑体字母 $\boldsymbol{a},\boldsymbol{b},\boldsymbol{r}$ 表示。

\begin{definition}
只考虑大小与方向、而不管起点位置的向量称为\textbf{自由向量}（本附录所讨论的向量均指自由向量）。若两个向量经平移后能完全重合，即它们的\textbf{模相等且方向相同}，则称这两个向量\textbf{相等}，记作 $\boldsymbol{a}=\boldsymbol{b}$。向量 $\boldsymbol{a}$ 的大小叫做它的\textbf{模}，记作 $|\boldsymbol{a}|$ 或 $|\overrightarrow{AB}|$。
\end{definition}

\begin{definition}
模等于 $1$ 的向量叫做\textbf{单位向量}；模等于 $0$ 的向量叫做\textbf{零向量}，记作 $\boldsymbol{0}$，规定零向量的方向任意。与 $\boldsymbol{a}$ 模相同、方向相反的向量叫做 $\boldsymbol{a}$ 的\textbf{负向量}，记作 $-\boldsymbol{a}$。
\end{definition}

\begin{definition}
设 $\boldsymbol{a},\boldsymbol{b}$ 是两个非零向量，任取空间一点 $O$，作 $\overrightarrow{OA}=\boldsymbol{a}$，$\overrightarrow{OB}=\boldsymbol{b}$，规定不超过 $\pi$ 的角 $\angle AOB$ 为 $\boldsymbol{a}$ 与 $\boldsymbol{b}$ 的\textbf{夹角}，记作 $(\boldsymbol{a},\boldsymbol{b})$ 或 $\langle\boldsymbol{a},\boldsymbol{b}\rangle$，于是 $0\le\langle\boldsymbol{a},\boldsymbol{b}\rangle\le\pi$。
\begin{enumerate}
\item 若 $\langle\boldsymbol{a},\boldsymbol{b}\rangle=0$ 或 $\pi$，称 $\boldsymbol{a}$ 与 $\boldsymbol{b}$ \textbf{平行}（共线），记作 $\boldsymbol{a}\parallel\boldsymbol{b}$；
\item 若 $\langle\boldsymbol{a},\boldsymbol{b}\rangle=\frac{\pi}{2}$，称 $\boldsymbol{a}$ 与 $\boldsymbol{b}$ \textbf{垂直}，记作 $\boldsymbol{a}\perp\boldsymbol{b}$。
\end{enumerate}
若 $\boldsymbol{a},\boldsymbol{b}$ 中有一个是零向量，规定它们的夹角可在 $[0,\pi]$ 内任意取值；因此零向量与任何向量都既平行又垂直。
\end{definition}

设有 $n\ (n\ge3)$ 个向量，把它们平移到同一起点，若各终点与这个公共起点在同一个平面上，就称这 $n$ 个向量\textbf{共面}。

\subsection{向量的加减法}

\begin{definition}[向量的加法]
设有两个向量 $\boldsymbol{a},\boldsymbol{b}$，任取一点 $A$，作 $\overrightarrow{AB}=\boldsymbol{a}$，再以 $B$ 为起点作 $\overrightarrow{BC}=\boldsymbol{b}$，则向量 $\overrightarrow{AC}=\boldsymbol{c}$ 叫做 $\boldsymbol{a}$ 与 $\boldsymbol{b}$ 的\textbf{和}，记作 $\boldsymbol{c}=\boldsymbol{a}+\boldsymbol{b}$。
\end{definition}

上述作法称为向量相加的\textbf{三角形法则}；当 $\boldsymbol{a},\boldsymbol{b}$ 不平行时，以 $\overrightarrow{AB}=\boldsymbol{a}$、$\overrightarrow{AD}=\boldsymbol{b}$ 为邻边作平行四边形，其对角线向量 $\overrightarrow{AC}=\boldsymbol{a}+\boldsymbol{b}$，这称为\textbf{平行四边形法则}（力学的合力法则）。

\begin{theorem}
向量的加法满足
\[
\boldsymbol{a}+\boldsymbol{b}=\boldsymbol{b}+\boldsymbol{a},\qquad
(\boldsymbol{a}+\boldsymbol{b})+\boldsymbol{c}=\boldsymbol{a}+(\boldsymbol{b}+\boldsymbol{c}).
\]
\end{theorem}

由交换律与结合律，$n$ 个向量相加可任意分组、任意换序，并有\textbf{多边形法则}：以前一个向量的终点作为次一个向量的起点相继作出这 $n$ 个向量，则以第一个向量的起点为起点、以最后一个向量的终点为终点的向量就是它们的和。

两个向量 $\boldsymbol{a}$ 与 $\boldsymbol{b}$ 的\textbf{差}规定为 $\boldsymbol{b}-\boldsymbol{a}=\boldsymbol{b}+(-\boldsymbol{a})$。把 $\boldsymbol{a},\boldsymbol{b}$ 平移到同一起点 $O$，则从 $\boldsymbol{a}$ 的终点指向 $\boldsymbol{b}$ 的终点的向量就是 $\boldsymbol{b}-\boldsymbol{a}$。

\keyformula{\big||\boldsymbol{a}|-|\boldsymbol{b}|\big|\ \le\ |\boldsymbol{a}\pm\boldsymbol{b}|\ \le\ |\boldsymbol{a}|+|\boldsymbol{b}|}

左端等号当且仅当 $\boldsymbol{a},\boldsymbol{b}$ 反向时成立，右端等号当且仅当 $\boldsymbol{a},\boldsymbol{b}$ 同向时成立。

\subsection{向量与数的乘法}

\begin{definition}
向量 $\boldsymbol{a}$ 与实数 $\lambda$ 的乘积 $\lambda\boldsymbol{a}$ 规定为一个向量：它的模 $|\lambda\boldsymbol{a}|=|\lambda|\,|\boldsymbol{a}|$；它的方向当 $\lambda>0$ 时与 $\boldsymbol{a}$ 相同，当 $\lambda<0$ 时与 $\boldsymbol{a}$ 相反，当 $\lambda=0$ 时 $\lambda\boldsymbol{a}=\boldsymbol{0}$。特别地 $1\boldsymbol{a}=\boldsymbol{a}$，$(-1)\boldsymbol{a}=-\boldsymbol{a}$。
\end{definition}

\begin{theorem}
数乘满足
\[
\lambda(\mu\boldsymbol{a})=(\lambda\mu)\boldsymbol{a},\qquad
(\lambda+\mu)\boldsymbol{a}=\lambda\boldsymbol{a}+\mu\boldsymbol{a},\qquad
\lambda(\boldsymbol{a}+\boldsymbol{b})=\lambda\boldsymbol{a}+\lambda\boldsymbol{b}.
\]
\end{theorem}

向量相加与数乘向量统称为\textbf{向量的线性运算}。取 $\lambda=1/|\boldsymbol{a}|$，得一个与 $\boldsymbol{a}$ 同方向、模为 $1$ 的向量，称为 $\boldsymbol{a}$ 的\textbf{单位化}：

\keyformula{\boldsymbol{e}_{\boldsymbol{a}}=\frac{\boldsymbol{a}}{|\boldsymbol{a}|}\qquad(\boldsymbol{a}\neq\boldsymbol{0})}

\begin{theorem}[两向量平行的充要条件]
设 $\boldsymbol{a}\neq\boldsymbol{0}$，则 $\boldsymbol{b}\parallel\boldsymbol{a}$ 的充分必要条件是：存在\textbf{唯一}的实数 $\lambda$，使 $\boldsymbol{b}=\lambda\boldsymbol{a}$。
\end{theorem}

必要性：若 $\boldsymbol{b}\parallel\boldsymbol{a}$，取 $\lambda=\pm\dfrac{|\boldsymbol{b}|}{|\boldsymbol{a}|}$（同向取正号、反向取负号）。唯一性：若 $\lambda\boldsymbol{a}=\mu\boldsymbol{a}$，则 $(\lambda-\mu)\boldsymbol{a}=\boldsymbol{0}$，而 $\boldsymbol{a}\neq\boldsymbol{0}$，故 $\lambda=\mu$。此定理是建立数轴的理论依据。

\subsection{空间直角坐标系与向量的坐标}

在空间取定一点 $O$ 和三个两两垂直的单位向量 $\boldsymbol{i},\boldsymbol{j},\boldsymbol{k}$，就确定了三条都以 $O$ 为原点的两两垂直的数轴，依次称为 $x$ 轴（横轴）、$y$ 轴（纵轴）、$z$ 轴（竖轴），统称\textbf{坐标轴}；它们构成\textbf{空间直角坐标系}，记作 $O$-$xyz$。当 $\boldsymbol{i},\boldsymbol{j},\boldsymbol{k}$ 的正向符合\textbf{右手规则}（右手四指从 $x$ 轴正向以 $\pi/2$ 转向 $y$ 轴正向时，大拇指的指向就是 $z$ 轴正向）时，称为右手坐标系。三条坐标轴中任意两条确定的平面统称\textbf{坐标面}，分别记作 $xOy$ 面、$yOz$ 面、$zOx$ 面；三个坐标面把空间分成八个部分，每一部分叫做一个\textbf{卦限}。

任给向量 $\boldsymbol{r}$，作 $\overrightarrow{OM}=\boldsymbol{r}$，以 $OM$ 为对角线、三条坐标轴为棱作长方体，则

\[
\boldsymbol{r}=\overrightarrow{OM}=\overrightarrow{OP}+\overrightarrow{OQ}+\overrightarrow{OR}
=x\boldsymbol{i}+y\boldsymbol{j}+z\boldsymbol{k}.
\]

上式称为向量 $\boldsymbol{r}$ 的\textbf{坐标分解式}，$x\boldsymbol{i},\,y\boldsymbol{j},\,z\boldsymbol{k}$ 称为 $\boldsymbol{r}$ 沿三个坐标轴方向的\textbf{分向量}，有序数 $x,y,z$ 称为 $\boldsymbol{r}$ 的\textbf{坐标}，记作 $\boldsymbol{r}=(x,y,z)$。

点 $M$ 与该点的向径 $\overrightarrow{OM}$ 有相同的坐标，故记号 $(x,y,z)$ 既可表示点 $M$，又可表示向量 $\overrightarrow{OM}$；但点是几何对象、向量是代数对象，\textbf{运算只对向量进行}。

坐标面与坐标轴上的点，其坐标有显著特征：$yOz$ 面上 $x=0$，$zOx$ 面上 $y=0$，$xOy$ 面上 $z=0$；$x$ 轴上 $y=z=0$，$y$ 轴上 $z=x=0$，$z$ 轴上 $x=y=0$；原点处 $x=y=z=0$。

\subsection{利用坐标作向量的线性运算}

设 $\boldsymbol{a}=(a_x,a_y,a_z)$，$\boldsymbol{b}=(b_x,b_y,b_z)$，$\lambda$ 为实数，则

\[
\boldsymbol{a}\pm\boldsymbol{b}=(a_x\pm b_x,\ a_y\pm b_y,\ a_z\pm b_z),\qquad
\lambda\boldsymbol{a}=(\lambda a_x,\ \lambda a_y,\ \lambda a_z).
\]

即对向量作加、减与数乘，只需对各个坐标分别作相应的数量运算。

设 $A(x_1,y_1,z_1)$、$B(x_2,y_2,z_2)$，则

\keyformula{\overrightarrow{AB}=(x_2-x_1,\ y_2-y_1,\ z_2-z_1)}

由平行的充要条件，当 $\boldsymbol{a}\neq\boldsymbol{0}$ 时

\keyformula{\boldsymbol{b}\parallel\boldsymbol{a}\iff\frac{b_x}{a_x}=\frac{b_y}{a_y}=\frac{b_z}{a_z}}

这里约定：某个分母为零时，相应的分子也必须为零。即\textbf{两向量平行等价于它们的坐标成比例}。

\textbf{分点坐标公式}\quad 设 $A(x_1,y_1,z_1)$、$B(x_2,y_2,z_2)$，点 $M$ 在直线 $AB$ 上且 $\overrightarrow{AM}=\lambda\overrightarrow{MB}$（$\lambda\neq-1$），则

\keyformula{\overrightarrow{OM}=\frac{\overrightarrow{OA}+\lambda\overrightarrow{OB}}{1+\lambda}}

即 $x=\dfrac{x_1+\lambda x_2}{1+\lambda}$，$y=\dfrac{y_1+\lambda y_2}{1+\lambda}$，$z=\dfrac{z_1+\lambda z_2}{1+\lambda}$。当 $\lambda=1$ 时得线段 $AB$ 的\textbf{中点}坐标 $\left(\dfrac{x_1+x_2}{2},\dfrac{y_1+y_2}{2},\dfrac{z_1+z_2}{2}\right)$。

\subsection{向量的模、方向角与投影}

由勾股定理，向量 $\boldsymbol{r}=(x,y,z)$ 的模

\keyformula{|\boldsymbol{r}|=\sqrt{x^{2}+y^{2}+z^{2}}}

从而点 $A(x_1,y_1,z_1)$ 与点 $B(x_2,y_2,z_2)$ 之间的距离为 $\overrightarrow{AB}$ 的模：

\keyformula{|AB|=\sqrt{(x_2-x_1)^{2}+(y_2-y_1)^{2}+(z_2-z_1)^{2}}}

\begin{definition}
非零向量 $\boldsymbol{r}$ 与三条坐标轴的夹角 $\alpha,\beta,\gamma$（$0\le\alpha,\beta,\gamma\le\pi$）称为 $\boldsymbol{r}$ 的\textbf{方向角}，$\cos\alpha,\cos\beta,\cos\gamma$ 称为 $\boldsymbol{r}$ 的\textbf{方向余弦}。
\end{definition}

\keyformula{\cos\alpha=\frac{x}{|\boldsymbol{r}|},\qquad
\cos\beta=\frac{y}{|\boldsymbol{r}|},\qquad
\cos\gamma=\frac{z}{|\boldsymbol{r}|}}

\keyformula{\cos^{2}\alpha+\cos^{2}\beta+\cos^{2}\gamma=1}

后一式表明：以方向余弦为坐标的向量正是与 $\boldsymbol{r}$ 同方向的单位向量，即
$\boldsymbol{e}_{\boldsymbol{r}}=(\cos\alpha,\ \cos\beta,\ \cos\gamma)$。

\begin{definition}
设数轴 $u$ 的单位向量为 $\boldsymbol{e}$，向量 $\boldsymbol{a}$ 与 $\boldsymbol{e}$ 的夹角为 $\varphi$，称 $|\boldsymbol{a}|\cos\varphi$ 为向量 $\boldsymbol{a}$ 在轴 $u$ 上的\textbf{投影}，记作 $\mathrm{Prj}_u\boldsymbol{a}$。
\end{definition}

投影具有线性性：

\[
\mathrm{Prj}_u(\boldsymbol{a}+\boldsymbol{b})=\mathrm{Prj}_u\boldsymbol{a}+\mathrm{Prj}_u\boldsymbol{b},\qquad
\mathrm{Prj}_u(\lambda\boldsymbol{a})=\lambda\,\mathrm{Prj}_u\boldsymbol{a}.
\]

向量的坐标就是它在各坐标轴上的投影：$a_x=\mathrm{Prj}_x\boldsymbol{a}$，$a_y=\mathrm{Prj}_y\boldsymbol{a}$，$a_z=\mathrm{Prj}_z\boldsymbol{a}$。投影是有符号的量，可正可负。

\corepoint{
（1）$\boldsymbol{a}\parallel\boldsymbol{b}$ 的坐标条件要写成“坐标成比例”，但\textbf{分母为零时必须单独讨论}：例如 $\boldsymbol{a}=(1,0,2)$、$\boldsymbol{b}=(2,0,4)$ 平行，此时 $\frac{0}{0}$ 无意义，只能写作 $b_x=2a_x,\ b_y=0,\ b_z=2a_z$。\\
（2）方向角的范围是 $[0,\pi]$，故 $\cos\alpha$ 可正可负；由 $\cos^{2}\alpha+\cos^{2}\beta+\cos^{2}\gamma=1$ 反求第三个方向余弦时，必须由点所在卦限（坐标符号）确定正负号，不能只取正根。\\
（3）记号 $(x,y,z)$ 兼指点与向量，须从上下文判断；对“点”不能作加减、数乘运算。
}

\section{数量积、向量积与混合积}

\subsection{两向量的数量积}

\begin{definition}
设 $\boldsymbol{a},\boldsymbol{b}$ 是两个向量，称 $|\boldsymbol{a}|\,|\boldsymbol{b}|\cos\langle\boldsymbol{a},\boldsymbol{b}\rangle$ 为 $\boldsymbol{a}$ 与 $\boldsymbol{b}$ 的\textbf{数量积}（内积、点积），记作 $\boldsymbol{a}\cdot\boldsymbol{b}$，即
\[
\boldsymbol{a}\cdot\boldsymbol{b}=|\boldsymbol{a}|\,|\boldsymbol{b}|\cos\langle\boldsymbol{a},\boldsymbol{b}\rangle .
\]
\end{definition}

物理背景：物体在恒力 $\boldsymbol{F}$ 作用下沿直线从 $M_1$ 移到 $M_2$，位移 $\boldsymbol{s}=\overrightarrow{M_1M_2}$，则力所做的功 $W=\boldsymbol{F}\cdot\boldsymbol{s}$。

设 $\boldsymbol{a}=(a_x,a_y,a_z)$，$\boldsymbol{b}=(b_x,b_y,b_z)$，由分配律与 $\boldsymbol{i}\cdot\boldsymbol{i}=\boldsymbol{j}\cdot\boldsymbol{j}=\boldsymbol{k}\cdot\boldsymbol{k}=1$、$\boldsymbol{i}\cdot\boldsymbol{j}=\boldsymbol{j}\cdot\boldsymbol{k}=\boldsymbol{k}\cdot\boldsymbol{i}=0$ 得坐标公式：

\keyformula{\boldsymbol{a}\cdot\boldsymbol{b}=a_xb_x+a_yb_y+a_zb_z}

\begin{property}
数量积满足下列运算律：
\begin{enumerate}
\item 交换律 $\boldsymbol{a}\cdot\boldsymbol{b}=\boldsymbol{b}\cdot\boldsymbol{a}$；
\item 分配律 $(\boldsymbol{a}+\boldsymbol{b})\cdot\boldsymbol{c}=\boldsymbol{a}\cdot\boldsymbol{c}+\boldsymbol{b}\cdot\boldsymbol{c}$；
\item 数乘结合律 $(\lambda\boldsymbol{a})\cdot\boldsymbol{b}=\lambda(\boldsymbol{a}\cdot\boldsymbol{b})$；
\item $\boldsymbol{a}\cdot\boldsymbol{a}=|\boldsymbol{a}|^{2}$，由此 $|\boldsymbol{a}|=\sqrt{\boldsymbol{a}\cdot\boldsymbol{a}}$。
\end{enumerate}
\end{property}

由定义与坐标公式即得两向量夹角的余弦：

\keyformula{\cos\langle\boldsymbol{a},\boldsymbol{b}\rangle=
\frac{\boldsymbol{a}\cdot\boldsymbol{b}}{|\boldsymbol{a}|\,|\boldsymbol{b}|}=
\frac{a_xb_x+a_yb_y+a_zb_z}
{\sqrt{a_x^{2}+a_y^{2}+a_z^{2}}\,\sqrt{b_x^{2}+b_y^{2}+b_z^{2}}}}

\begin{theorem}[垂直的充要条件]
设 $\boldsymbol{a},\boldsymbol{b}$ 为非零向量，则 $\boldsymbol{a}\perp\boldsymbol{b}\iff\boldsymbol{a}\cdot\boldsymbol{b}=0$。
\end{theorem}

数量积与投影的关系是：$\boldsymbol{a}\cdot\boldsymbol{b}=|\boldsymbol{a}|\,\mathrm{Prj}_{\boldsymbol{a}}\boldsymbol{b}=|\boldsymbol{b}|\,\mathrm{Prj}_{\boldsymbol{b}}\boldsymbol{a}$；特别当 $\boldsymbol{e}$ 为单位向量时 $\mathrm{Prj}_{\boldsymbol{e}}\boldsymbol{a}=\boldsymbol{a}\cdot\boldsymbol{e}$。

几何意义：数量积等于其中一个向量的模与另一个向量在这个向量方向上的投影的乘积。

\subsection{两向量的向量积}

\begin{definition}
对空间中任意两个向量 $\boldsymbol{a}$ 与 $\boldsymbol{b}$，规定它们的\textbf{向量积}（叉积、外积）$\boldsymbol{a}\times\boldsymbol{b}$ 如下：
\begin{enumerate}
\item 若 $\boldsymbol{a},\boldsymbol{b}$ 不平行，则 $\boldsymbol{a}\times\boldsymbol{b}$ 是一个向量，满足\\
\quad（i）$|\boldsymbol{a}\times\boldsymbol{b}|=|\boldsymbol{a}|\,|\boldsymbol{b}|\sin\langle\boldsymbol{a},\boldsymbol{b}\rangle$；\\
\quad（ii）$\boldsymbol{a}\times\boldsymbol{b}$ 垂直于 $\boldsymbol{a}$ 与 $\boldsymbol{b}$，且 $\boldsymbol{a},\boldsymbol{b},\boldsymbol{a}\times\boldsymbol{b}$ 依次构成\textbf{右手系}（右手四指从 $\boldsymbol{a}$ 转向 $\boldsymbol{b}$ 时，大拇指指向 $\boldsymbol{a}\times\boldsymbol{b}$ 的方向）；
\item 若 $\boldsymbol{a}\parallel\boldsymbol{b}$，规定 $\boldsymbol{a}\times\boldsymbol{b}=\boldsymbol{0}$。
\end{enumerate}
注意其中 $\boldsymbol{a}$ 与 $\boldsymbol{b}$ 的次序不能颠倒。
\end{definition}

坐标公式可按“$\boldsymbol{i},\boldsymbol{j},\boldsymbol{k}$ 作第一行、$\boldsymbol{a}$ 的坐标作第二行、$\boldsymbol{b}$ 的坐标作第三行”的三阶行列式展开：

\keyformula{\boldsymbol{a}\times\boldsymbol{b}=
\begin{vmatrix}
\boldsymbol{i}&\boldsymbol{j}&\boldsymbol{k}\\
a_x&a_y&a_z\\
b_x&b_y&b_z
\end{vmatrix}
=\begin{vmatrix}
a_y&a_z\\ b_y&b_z
\end{vmatrix}\boldsymbol{i}
-\begin{vmatrix}
a_x&a_z\\ b_x&b_z
\end{vmatrix}\boldsymbol{j}
+\begin{vmatrix}
a_x&a_y\\ b_x&b_y
\end{vmatrix}\boldsymbol{k}}

即 $\boldsymbol{a}\times\boldsymbol{b}=(a_yb_z-a_zb_y,\ a_zb_x-a_xb_z,\ a_xb_y-a_yb_x)$。展开时第二项带负号，这是最容易出错的地方。

\begin{property}
向量积满足：
\begin{enumerate}
\item 反交换律 $\boldsymbol{a}\times\boldsymbol{b}=-\boldsymbol{b}\times\boldsymbol{a}$，特别地 $\boldsymbol{a}\times\boldsymbol{a}=\boldsymbol{0}$；
\item $\boldsymbol{a}\times\boldsymbol{b}\perp\boldsymbol{a}$，$\boldsymbol{a}\times\boldsymbol{b}\perp\boldsymbol{b}$；
\item 分配律 $(\boldsymbol{a}+\boldsymbol{b})\times\boldsymbol{c}=\boldsymbol{a}\times\boldsymbol{c}+\boldsymbol{b}\times\boldsymbol{c}$；
\item 数乘结合 $(\lambda\boldsymbol{a})\times\boldsymbol{b}=\lambda(\boldsymbol{a}\times\boldsymbol{b})=\boldsymbol{a}\times(\lambda\boldsymbol{b})$；
\item $\boldsymbol{a}\parallel\boldsymbol{b}\iff\boldsymbol{a}\times\boldsymbol{b}=\boldsymbol{0}$（$\boldsymbol{a},\boldsymbol{b}$ 为非零向量）。
\end{enumerate}
\end{property}

几何意义：$|\boldsymbol{a}\times\boldsymbol{b}|$ 等于以 $\boldsymbol{a},\boldsymbol{b}$ 为邻边的平行四边形的面积；以 $\boldsymbol{a},\boldsymbol{b}$ 为邻边的三角形面积为 $\frac12|\boldsymbol{a}\times\boldsymbol{b}|$。物理上力矩 $\boldsymbol{M}=\boldsymbol{r}\times\boldsymbol{F}$，刚体绕定轴以角速度 $\boldsymbol{\omega}$ 转动时，点 $M$（向径 $\boldsymbol{r}$）的线速度 $\boldsymbol{v}=\boldsymbol{\omega}\times\boldsymbol{r}$。

\subsection{向量的混合积}

\begin{definition}
已知三个向量 $\boldsymbol{a},\boldsymbol{b},\boldsymbol{c}$，先作 $\boldsymbol{a}$ 与 $\boldsymbol{b}$ 的向量积 $\boldsymbol{a}\times\boldsymbol{b}$，再与 $\boldsymbol{c}$ 作数量积 $(\boldsymbol{a}\times\boldsymbol{b})\cdot\boldsymbol{c}$，所得的数量叫做这三个向量的\textbf{混合积}，记作 $[\boldsymbol{a}\boldsymbol{b}\boldsymbol{c}]$。
\end{definition}

\keyformula{[\boldsymbol{a}\boldsymbol{b}\boldsymbol{c}]=(\boldsymbol{a}\times\boldsymbol{b})\cdot\boldsymbol{c}=
\begin{vmatrix}
a_x&a_y&a_z\\
b_x&b_y&b_z\\
c_x&c_y&c_z
\end{vmatrix}}

\begin{property}
混合积满足：
\begin{enumerate}
\item \textbf{轮换对称性} $[\boldsymbol{a}\boldsymbol{b}\boldsymbol{c}]=[\boldsymbol{b}\boldsymbol{c}\boldsymbol{a}]=[\boldsymbol{c}\boldsymbol{a}\boldsymbol{b}]$；
\item 任意对换两个向量，混合积变号：$[\boldsymbol{a}\boldsymbol{b}\boldsymbol{c}]=-[\boldsymbol{b}\boldsymbol{a}\boldsymbol{c}]$；
\item $[\boldsymbol{a}\boldsymbol{b}\boldsymbol{c}]=0$ 的充要条件是 $\boldsymbol{a},\boldsymbol{b},\boldsymbol{c}$ \textbf{共面}（含其中有零向量或有两者平行的情况）。
\end{enumerate}
\end{property}

几何意义：$|[\boldsymbol{a}\boldsymbol{b}\boldsymbol{c}]|$ 等于以 $\boldsymbol{a},\boldsymbol{b},\boldsymbol{c}$ 为棱的平行六面体的体积；以它们为棱的四面体体积为 $\frac16|[\boldsymbol{a}\boldsymbol{b}\boldsymbol{c}]|$。混合积的\textbf{符号}由 $\boldsymbol{a},\boldsymbol{b},\boldsymbol{c}$ 构成右手系（为正）还是左手系（为负）决定。

\corepoint{
（1）数量积的结果是\textbf{数}，向量积的结果是\textbf{向量}，混合积的结果是\textbf{数}；写 $\boldsymbol{a}\cdot\boldsymbol{b}\cdot\boldsymbol{c}$ 没有意义，必须写成 $(\boldsymbol{a}\cdot\boldsymbol{b})\boldsymbol{c}$ 或 $[\boldsymbol{a}\boldsymbol{b}\boldsymbol{c}]$。\\
（2）向量积\textbf{不满足结合律}：一般地 $(\boldsymbol{a}\times\boldsymbol{b})\times\boldsymbol{c}\neq\boldsymbol{a}\times(\boldsymbol{b}\times\boldsymbol{c})$，只有分配律可用。\\
（3）向量积坐标公式的第二个分量带负号，务必按行列式展开，不要凭记忆写 $(a_yb_z-a_zb_y,\ a_zb_x-a_xb_z,\ \dots)$ 时把符号弄反。\\
（4）用混合积判断共面、求体积时，结论是“混合积为零”与“取绝对值”，符号本身只表示左右手系。
}

\section{平面及其方程}

\subsection{平面的点法式方程}

\begin{definition}
垂直于平面 $\pi$ 的非零向量 $\boldsymbol{n}=(A,B,C)$ 称为平面 $\pi$ 的\textbf{法向量}。
\end{definition}

设平面 $\pi$ 过点 $M_0(x_0,y_0,z_0)$，且以 $\boldsymbol{n}=(A,B,C)\neq\boldsymbol{0}$ 为法向量。点 $M(x,y,z)$ 在 $\pi$ 上的充要条件是 $\overrightarrow{M_0M}\perp\boldsymbol{n}$，即 $\boldsymbol{n}\cdot\overrightarrow{M_0M}=0$，于是

\keyformula{A(x-x_0)+B(y-y_0)+C(z-z_0)=0}

称其为平面 $\pi$ 的\textbf{点法式方程}。几何上“一点 + 一个法向量”唯一确定一张平面。

\subsection{平面的一般方程}

把点法式展开，得

\[
Ax+By+Cz+D=0\qquad(A,B,C\ \text{不全为零}),
\]

称为平面的\textbf{一般方程}；其法向量就是 $\boldsymbol{n}=(A,B,C)$。反之，任何一个三元一次方程都表示一张平面。三元一次方程中三个系数 $A,B,C$ 的取值决定了平面的位置，常见特殊情形如下。

\begin{table}[htbp]
\centering
\renewcommand{\tablename}{表}
\caption{平面一般方程 $Ax+By+Cz+D=0$ 的特殊位置}
\begin{tabular}{@{}ll@{}}
\toprule
方程特征 & 图形 \\
\midrule
$D=0$ & 过原点 \\
$A=0$ & 平行于 $x$ 轴 \\
$A=D=0$ & 通过 $x$ 轴 \\
$A=B=0$ & 平行于 $xOy$ 面（即 $z=-\frac{D}{C}$）\\
$A=B=D=0$ & 就是 $xOy$ 面（$z=0$）\\
\bottomrule
\end{tabular}
\end{table}

若平面与三条坐标轴分别交于 $P(a,0,0)$、$Q(0,b,0)$、$R(0,0,c)$，且 $abc\neq0$，则平面方程可写成

\keyformula{\frac{x}{a}+\frac{y}{b}+\frac{z}{c}=1}

称为平面的\textbf{截距式方程}，$a,b,c$ 依次称为平面在 $x$ 轴、$y$ 轴、$z$ 轴上的\textbf{截距}。截距式便于快速作图，但过原点或平行于某坐标轴的平面没有截距式。

\textbf{三点式}\quad 过不共线三点 $M_1,M_2,M_3$ 的平面，取法向量 $\boldsymbol{n}=\overrightarrow{M_1M_2}\times\overrightarrow{M_1M_3}$，再按点法式写出即可；三点共线的判据是 $\overrightarrow{M_1M_2}\times\overrightarrow{M_1M_3}=\boldsymbol{0}$。

\subsection{两平面的夹角}

\begin{definition}
两平面的\textbf{夹角}规定为它们法向量的夹角（取锐角或直角），即 $\theta\in\left[0,\frac{\pi}{2}\right]$。
\end{definition}

设 $\pi_1:A_1x+B_1y+C_1z+D_1=0$，$\pi_2:A_2x+B_2y+C_2z+D_2=0$，法向量分别为 $\boldsymbol{n}_1,\boldsymbol{n}_2$，则

\keyformula{\cos\theta=\frac{|\boldsymbol{n}_1\cdot\boldsymbol{n}_2|}{|\boldsymbol{n}_1|\,|\boldsymbol{n}_2|}=
\frac{|A_1A_2+B_1B_2+C_1C_2|}
{\sqrt{A_1^{2}+B_1^{2}+C_1^{2}}\,\sqrt{A_2^{2}+B_2^{2}+C_2^{2}}}}

两平面的位置关系由法向量与常数项共同决定：

\begin{table}[htbp]
\centering
\renewcommand{\tablename}{表}
\caption{两平面位置关系的判据}
\begin{tabular}{@{}ll@{}}
\toprule
位置关系 & 充要条件 \\
\midrule
$\pi_1\perp\pi_2$ & $A_1A_2+B_1B_2+C_1C_2=0$ \\
$\pi_1\parallel\pi_2$（不重合） & $\dfrac{A_1}{A_2}=\dfrac{B_1}{B_2}=\dfrac{C_1}{C_2}\neq\dfrac{D_1}{D_2}$ \\
$\pi_1$ 与 $\pi_2$ 重合 & $\dfrac{A_1}{A_2}=\dfrac{B_1}{B_2}=\dfrac{C_1}{C_2}=\dfrac{D_1}{D_2}$ \\
\bottomrule
\end{tabular}
\end{table}

\subsection{点到平面的距离}

设 $P_0(x_0,y_0,z_0)$ 是平面 $\pi:Ax+By+Cz+D=0$ 外一点，在 $\pi$ 上任取一点 $P_1$，则 $\overrightarrow{P_1P_0}$ 在法向量 $\boldsymbol{n}$ 上的投影的绝对值就是所求距离：

\keyformula{d=\frac{|Ax_0+By_0+Cz_0+D|}{\sqrt{A^{2}+B^{2}+C^{2}}}}

推导要点：$d=\left|\mathrm{Prj}_{\boldsymbol{n}}\overrightarrow{P_1P_0}\right|=\dfrac{|\boldsymbol{n}\cdot\overrightarrow{P_1P_0}|}{|\boldsymbol{n}|}$，而 $\boldsymbol{n}\cdot\overrightarrow{P_1P_0}=A(x_0-x_1)+B(y_0-y_1)+C(z_0-z_1)=Ax_0+By_0+Cz_0+D$（用了 $P_1$ 在平面上）。

由此可得两\textbf{平行}平面 $Ax+By+Cz+D_1=0$ 与 $Ax+By+Cz+D_2=0$ 之间的距离

\keyformula{d=\frac{|D_1-D_2|}{\sqrt{A^{2}+B^{2}+C^{2}}}}

\corepoint{
（1）点到平面距离公式中分子要加\textbf{绝对值}，否则得到的是有符号的“离差”，不能当距离。\\
（2）用平行平面距离公式前，必须先把两方程化成 $x,y,z$ 系数\textbf{完全相同}的形式：若给的是 $2x+4y-2z+1=0$ 与 $x+2y-z+3=0$，须先把后者乘 $2$（得 $2x+4y-2z+6=0$）再套公式。\\
（3）平面一般方程的三个系数不能同时为零；若只有一个系数为零（如 $A=0$），平面仍存在，只是平行于相应的坐标轴，不要误判为“退化”。
}

\section{空间直线及其方程}

\subsection{空间直线的一般方程}

空间直线可看作两张不平行的平面的交线。设

\[
\pi_1:A_1x+B_1y+C_1z+D_1=0,\qquad
\pi_2:A_2x+B_2y+C_2z+D_2=0,
\]

且 $\boldsymbol{n}_1=(A_1,B_1,C_1)$ 与 $\boldsymbol{n}_2=(A_2,B_2,C_2)$ 不平行（即 $A_1:B_1:C_1\neq A_2:B_2:C_2$），则方程组

\[
\begin{cases}
A_1x+B_1y+C_1z+D_1=0,\\
A_2x+B_2y+C_2z+D_2=0
\end{cases}
\]

表示一条空间直线，称为空间直线的\textbf{一般方程}。同一条直线可写成无数种一般方程（取不同的两张过它的平面），故一般方程不唯一。

\subsection{对称式方程与参数方程}

\begin{definition}
平行于直线 $L$ 的非零向量 $\boldsymbol{s}=(m,n,p)$ 称为 $L$ 的\textbf{方向向量}，$m,n,p$ 称为 $L$ 的\textbf{方向数}。
\end{definition}

设直线 $L$ 过点 $M_0(x_0,y_0,z_0)$，方向向量为 $\boldsymbol{s}=(m,n,p)$，则点 $M(x,y,z)$ 在 $L$ 上的充要条件是 $\overrightarrow{M_0M}\parallel\boldsymbol{s}$，即存在 $t$ 使 $\overrightarrow{M_0M}=t\boldsymbol{s}$。消去 $t$ 得

\keyformula{\frac{x-x_0}{m}=\frac{y-y_0}{n}=\frac{z-z_0}{p}}

称为直线的\textbf{对称式方程}（点向式方程）。约定：某个分母为零时，相应的分子也为零，例如 $m=0$ 时该式应理解为 $x=x_0$ 与 $\dfrac{y-y_0}{n}=\dfrac{z-z_0}{p}$ 联立。

保留参数 $t$ 即得\textbf{参数方程}

\keyformula{\begin{cases}
x=x_0+mt,\\
y=y_0+nt,\\
z=z_0+pt
\end{cases}\qquad(-\infty<t<+\infty)}

过两点 $M_1(x_1,y_1,z_1)$、$M_2(x_2,y_2,z_2)$ 的直线，取 $\boldsymbol{s}=\overrightarrow{M_1M_2}$，得\textbf{两点式}

\keyformula{\frac{x-x_1}{x_2-x_1}=\frac{y-y_1}{y_2-y_1}=\frac{z-z_1}{z_2-z_1}}

\subsection{由一般方程求对称式的步骤}

\begin{enumerate}
\item 求方向向量：$\boldsymbol{s}=\boldsymbol{n}_1\times\boldsymbol{n}_2$，它同时垂直于两个平面的法向量，故平行于交线；
\item 求直线上一点：在一般方程中令某个坐标（通常取使计算最简的那个，如 $z=0$）为常数，解出其余两个坐标，得 $M_0$；
\item 写出对称式；若 $\boldsymbol{s}$ 的某个分量为零，按约定写成“分子为零”的形式，不要写成除以零。
\end{enumerate}

\section{位置关系与距离}

\subsection{两直线的位置关系}

设
\[
L_1:\frac{x-x_1}{m_1}=\frac{y-y_1}{n_1}=\frac{z-z_1}{p_1},\qquad
L_2:\frac{x-x_2}{m_2}=\frac{y-y_2}{n_2}=\frac{z-z_2}{p_2},
\]
方向向量分别为 $\boldsymbol{s}_1=(m_1,n_1,p_1)$、$\boldsymbol{s}_2=(m_2,n_2,p_2)$，$M_1\in L_1$，$M_2\in L_2$。两直线的位置关系由 $\boldsymbol{s}_1,\boldsymbol{s}_2,\overrightarrow{M_1M_2}$ 三者的关系完全确定：

\begin{enumerate}
\item $L_1$ 与 $L_2$ \textbf{重合}：$\boldsymbol{s}_1\parallel\boldsymbol{s}_2$ 且 $M_1\in L_2$（即 $\overrightarrow{M_1M_2}\parallel\boldsymbol{s}_2$）；
\item $L_1\parallel L_2$（不重合）：$\boldsymbol{s}_1\parallel\boldsymbol{s}_2$ 但 $\overrightarrow{M_1M_2}$ 不平行于 $\boldsymbol{s}_2$；
\item $L_1$ 与 $L_2$ \textbf{相交}于一点：$\boldsymbol{s}_1,\boldsymbol{s}_2$ 不平行，且混合积 $[\boldsymbol{s}_1\boldsymbol{s}_2\overrightarrow{M_1M_2}]=0$（三向量共面）；
\item $L_1$ 与 $L_2$ \textbf{异面}：混合积 $[\boldsymbol{s}_1\boldsymbol{s}_2\overrightarrow{M_1M_2}]\neq0$（三向量不共面）。
\end{enumerate}

即：\textbf{先看方向向量是否平行，再用混合积是否为零区分“相交”与“异面”}。

\subsection{两直线的夹角}

两直线的夹角 $\theta$ 规定为它们方向向量的夹角，取锐角或直角，$\theta\in\left[0,\frac{\pi}{2}\right]$：

\keyformula{\cos\theta=\frac{|\boldsymbol{s}_1\cdot\boldsymbol{s}_2|}{|\boldsymbol{s}_1|\,|\boldsymbol{s}_2|}=
\frac{|m_1m_2+n_1n_2+p_1p_2|}
{\sqrt{m_1^{2}+n_1^{2}+p_1^{2}}\,\sqrt{m_2^{2}+n_2^{2}+p_2^{2}}}}

由此：$L_1\perp L_2\iff m_1m_2+n_1n_2+p_1p_2=0$；$L_1\parallel L_2\iff\dfrac{m_1}{m_2}=\dfrac{n_1}{n_2}=\dfrac{p_1}{p_2}$（分母为零时按约定理解）。

\subsection{直线与平面的位置关系}

设直线 $L:\dfrac{x-x_0}{m}=\dfrac{y-y_0}{n}=\dfrac{z-z_0}{p}$，平面 $\pi:Ax+By+Cz+D=0$，方向向量 $\boldsymbol{s}=(m,n,p)$，法向量 $\boldsymbol{n}=(A,B,C)$。

\begin{enumerate}
\item $L$ \textbf{在平面} $\pi$ 上：$\boldsymbol{s}\cdot\boldsymbol{n}=0$ 且 $Ax_0+By_0+Cz_0+D=0$；
\item $L\parallel\pi$（不落在平面上）：$\boldsymbol{s}\cdot\boldsymbol{n}=0$ 而 $Ax_0+By_0+Cz_0+D\neq0$；
\item $L$ 与 $\pi$ \textbf{相交}于一点：$\boldsymbol{s}\cdot\boldsymbol{n}\neq0$，交点的求法是把 $L$ 的参数方程代入 $\pi$ 的方程解出 $t$，再回代；
\item $L\perp\pi$：$\boldsymbol{s}\parallel\boldsymbol{n}$，即 $\dfrac{m}{A}=\dfrac{n}{B}=\dfrac{p}{C}$。
\end{enumerate}

\begin{definition}
直线 $L$ 与它在平面 $\pi$ 上的\textbf{投影}所成的夹角 $\varphi$（$0\le\varphi\le\frac{\pi}{2}$）称为直线与平面的\textbf{夹角}。
\end{definition}

由于 $\boldsymbol{s}$ 与 $\boldsymbol{n}$ 的夹角与 $\varphi$ 互余（或为其补角），故用\textbf{正弦}：

\keyformula{\sin\varphi=\frac{|\boldsymbol{s}\cdot\boldsymbol{n}|}{|\boldsymbol{s}|\,|\boldsymbol{n}|}=
\frac{|Am+Bn+Cp|}
{\sqrt{A^{2}+B^{2}+C^{2}}\,\sqrt{m^{2}+n^{2}+p^{2}}}}

\corepoint{
（1）直线与平面夹角用 $\sin$，两直线夹角与两平面夹角用 $\cos$，三者都取绝对值使夹角落在 $\left[0,\frac{\pi}{2}\right]$；弄混是最高频的错误。\\
（2）“$\boldsymbol{s}\cdot\boldsymbol{n}=0$”只能说明直线平行于平面\textbf{或落在平面上}，必须再代一个点检验，才能区分这两种情形。\\
（3）两直线“方向向量不平行 + 混合积为零”才是相交；只验证混合积为零而忽略平行情形，会把重合或平行的直线误判为相交。
}

\subsection{点到直线的距离}

设点 $P(x_0,y_0,z_0)$ 与直线 $L:\dfrac{x-x_1}{m}=\dfrac{y-y_1}{n}=\dfrac{z-z_1}{p}$（$M_1(x_1,y_1,z_1)\in L$，方向向量 $\boldsymbol{s}=(m,n,p)$）。以 $\overrightarrow{M_1P}$ 与 $\boldsymbol{s}$ 为邻边的平行四边形，其面积为 $|\boldsymbol{s}\times\overrightarrow{M_1P}|$，而该面积又等于 $|\boldsymbol{s}|\cdot d$，故

\keyformula{d=\frac{|\boldsymbol{s}\times\overrightarrow{M_1P}|}{|\boldsymbol{s}|}}

若 $L$ 以一般方程给出，先用 $\boldsymbol{s}=\boldsymbol{n}_1\times\boldsymbol{n}_2$ 求方向向量，并在 $L$ 上任取一点作 $M_1$。

\subsection{两平行直线与两异面直线间的距离}

\textbf{两平行直线间的距离}\quad 设 $L_1\parallel L_2$，公共方向向量 $\boldsymbol{s}$，$M_1\in L_1$，$M_2\in L_2$，则

\keyformula{d=\frac{|\boldsymbol{s}\times\overrightarrow{M_1M_2}|}{|\boldsymbol{s}|}}

与点到直线距离公式同形：在 $L_2$ 上任取一点，化为“点到直线”的情形。

\textbf{两异面直线间的距离}\quad 设
\[
L_1:\frac{x-x_1}{m_1}=\frac{y-y_1}{n_1}=\frac{z-z_1}{p_1},\qquad
L_2:\frac{x-x_2}{m_2}=\frac{y-y_2}{n_2}=\frac{z-z_2}{p_2},
\]
$\boldsymbol{s}_1=(m_1,n_1,p_1)$，$\boldsymbol{s}_2=(m_2,n_2,p_2)$，$M_1(x_1,y_1,z_1)$，$M_2(x_2,y_2,z_2)$。两异面直线的公垂线方向为 $\boldsymbol{s}_1\times\boldsymbol{s}_2$，$\overrightarrow{M_1M_2}$ 在该方向上的投影的绝对值就是 $L_1$ 与 $L_2$ 之间的距离：

\keyformula{d=\frac{|[\boldsymbol{s}_1\boldsymbol{s}_2\overrightarrow{M_1M_2}]|}{|\boldsymbol{s}_1\times\boldsymbol{s}_2|}}

注意：此式要求 $\boldsymbol{s}_1\times\boldsymbol{s}_2\neq\boldsymbol{0}$（即两直线不平行）；若两直线平行，公垂线方向不唯一，应改用上一式。

\corepoint{
（1）求距离的三条公式本质相同，都是“面积（或体积）除以底”：点到直线用平行四边形面积 $|\boldsymbol{s}\times\overrightarrow{M_1P}|$，异面直线用平行六面体体积 $|[\boldsymbol{s}_1\boldsymbol{s}_2\overrightarrow{M_1M_2}]|$。记住这个几何来源，就不必死记公式。\\
（2）异面直线距离公式的分母是 $|\boldsymbol{s}_1\times\boldsymbol{s}_2|$ 而不是 $|\boldsymbol{s}_1|\,|\boldsymbol{s}_2|$，两者相差 $\sin\langle\boldsymbol{s}_1,\boldsymbol{s}_2\rangle$。\\
（3）所有距离公式最后都要取绝对值；用混合积算体积时若得负数，只说明三向量成左手系，距离仍取绝对值。
}

\section{平面束}

\subsection{平面束的方程}

通过给定直线 $L$ 的所有平面的全体，称为通过 $L$ 的\textbf{平面束}。设 $L$ 由一般方程

\[
\pi_1:A_1x+B_1y+C_1z+D_1=0,\qquad
\pi_2:A_2x+B_2y+C_2z+D_2=0
\]

给出，且 $\boldsymbol{n}_1=(A_1,B_1,C_1)$ 与 $\boldsymbol{n}_2=(A_2,B_2,C_2)$ 不平行（即 $\pi_1,\pi_2$ 确实交于一条直线）。对任意实数 $\lambda$，方程

\[
A_1x+B_1y+C_1z+D_1+\lambda\bigl(A_2x+B_2y+C_2z+D_2\bigr)=0
\]

仍是 $x,y,z$ 的一次方程，因而表示一张平面；又因为 $L$ 上每点同时满足两个括号为零，所以 $L$ 上的点都满足它。于是它表示过 $L$ 的平面。反之，除 $\pi_2$ 以外的任一过 $L$ 的平面都可由某个 $\lambda$ 表示，故

\keyformula{\pi_1+\lambda\pi_2=0\quad(\lambda\in\mathbb{R})}

称为过直线 $L$ 的\textbf{平面束方程}（单参数形式）。若需要把 $\pi_2$ 本身也包含进来，改用两个参数 $\lambda,\mu$ 不同时为零的齐次组合

\keyformula{\lambda\bigl(A_1x+B_1y+C_1z+D_1\bigr)+\mu\bigl(A_2x+B_2y+C_2z+D_2\bigr)=0}

此时表示过 $L$ 的\textbf{全部}平面。

\subsection{用平面束求平面方程的步骤}

\begin{enumerate}
\item 写出平面束方程 $\pi_1+\lambda\pi_2=0$，其中 $\lambda$ 待定；
\item 把题设条件翻译成关于 $\lambda$ 的方程：过某点 $\Rightarrow$ 代入该点坐标为零；垂直于某平面 $\Rightarrow$ 两法向量数量积为零；平行于某向量 $\Rightarrow$ 法向量与该向量数量积为零；与某平面成定角 $\Rightarrow$ 用法向量夹角余弦公式；
\item 解出 $\lambda$ 后代回，化简即得所求平面；
\item 若步骤 2 得到的方程无解或恒成立，说明所求平面正是 $\pi_2$，需单独取出 $\pi_2$ 检验。
\end{enumerate}

\corepoint{
（1）单参数形式 $\pi_1+\lambda\pi_2=0$ \textbf{永远漏掉} $\pi_2$ 这一张平面；若所求平面的条件恰好只被 $\pi_2$ 满足，用单参数式会“无解”，此时应写双参数形式或直接验证 $\pi_2$。\\
（2）平面束只对\textbf{相交}的两平面成立。若 $\boldsymbol{n}_1\parallel\boldsymbol{n}_2$，则 $\pi_1+\lambda\pi_2=0$ 表示与这两平面\textbf{平行}的一族平面，不是过同一直线的平面束，判据是有没有“过同一条直线”。\\
（3）$\pi_1,\pi_2$ 的取法不唯一（同一条直线可用不同两张平面表示），取计算最简便的一组，例如让其中一个平面的某个系数为零。
}

\section{常见曲面}

\subsection{曲面方程的概念}

\begin{definition}
若曲面 $S$ 上每一点的坐标都满足方程 $F(x,y,z)=0$，而不在 $S$ 上的点的坐标都不满足这个方程，则称 $F(x,y,z)=0$ 为曲面 $S$ 的\textbf{方程}，$S$ 为该方程的\textbf{图形}。
\end{definition}

曲面研究有两个基本问题：\textbf{已知曲面求方程}——用轨迹法，把几何条件“翻译”成坐标的等式；\textbf{已知方程求图形}——用截痕法，靠截线形状推断曲面形状。空间曲线则看作两张曲面的交线，其一般方程为两个方程联立的方程组。

\subsection{旋转曲面}

\begin{definition}
一条平面曲线 $C$ 绕它所在平面内的一条定直线 $L$ 旋转一周所成的曲面，叫做\textbf{旋转曲面}；$C$ 称为\textbf{母线}，$L$ 称为\textbf{旋转轴}。
\end{definition}

设母线 $C$ 在 $yOz$ 面上，方程为 $f(y,z)=0$。把 $C$ 绕 $z$ 轴旋转，则旋转过程中 $z$ 坐标不变，而点到 $z$ 轴的距离 $\sqrt{x^{2}+y^{2}}$ 保持等于原来的 $|y|$，故旋转曲面的方程为

\keyformula{f\bigl(\pm\sqrt{x^{2}+y^{2}},\ z\bigr)=0}

同理，$yOz$ 面上的曲线 $f(y,z)=0$ 绕 $y$ 轴旋转所得曲面为 $f\bigl(y,\ \pm\sqrt{x^{2}+z^{2}}\bigr)=0$。

记忆口诀：\textbf{与旋转轴同名的变量不动，另一个变量换成另外两个变量平方和的平方根}（正负号成对出现，平方后往往自动消去）。例如 $yOz$ 面上的抛物线 $y^{2}=2pz$ 绕 $z$ 轴旋转，得旋转抛物面 $x^{2}+y^{2}=2pz$；$yOz$ 面上的椭圆 $\dfrac{y^{2}}{b^{2}}+\dfrac{z^{2}}{c^{2}}=1$ 绕 $z$ 轴旋转，得旋转椭球面 $\dfrac{x^{2}+y^{2}}{b^{2}}+\dfrac{z^{2}}{c^{2}}=1$。

\subsection{柱面}

\begin{definition}
平行于定直线并沿定曲线 $C$ 移动的直线 $L$ 所形成的曲面叫做\textbf{柱面}，$C$ 称为\textbf{准线}，$L$ 称为\textbf{母线}。
\end{definition}

\textbf{判据}\quad 若曲面方程中\emph{只含两个变量}（例如只含 $x,y$），则它表示柱面，母线平行于\emph{方程中不出现的那个变量}所对应的坐标轴（此例中母线平行于 $z$ 轴），该柱面在 $xOy$ 面上的截痕就是它的准线。

常见柱面：$\dfrac{x^{2}}{a^{2}}+\dfrac{y^{2}}{b^{2}}=1$ 为椭圆柱面；$\dfrac{x^{2}}{a^{2}}-\dfrac{y^{2}}{b^{2}}=1$ 为双曲柱面；$y^{2}=2px$ 为抛物柱面；$x^{2}+y^{2}=R^{2}$ 为圆柱面。它们的母线都平行于 $z$ 轴。

\subsection{二次曲面与截痕法}

\begin{definition}
三元二次方程所表示的曲面叫做\textbf{二次曲面}。通过适当的坐标变换（平移、旋转），二次曲面可化为若干种\textbf{标准方程}之一。
\end{definition}

\textbf{截痕法}是识别二次曲面类型的基本方法：

\begin{enumerate}
\item 判断对称性：方程中 $x,y,z$ 都以平方项（或一次项）出现时，曲面关于相应的坐标面、坐标轴、原点对称；
\item 求与三个坐标面的截痕：分别令 $x=0$、$y=0$、$z=0$，看所得平面曲线的类型；
\item 用平行平面截割：令 $z=h$（或 $x=h$、$y=h$），观察截线类型随 $h$ 的变化（椭圆的大小、双曲线的实轴方向、是否退化为直线或点）；
\item 综合对称性、截线类型与变化趋势，定出曲面名称与大致图形（是否有界、是否分叶、是否有直纹）。
\end{enumerate}

\subsection{二次曲面标准方程对照}

\begingroup
\renewcommand{\tablename}{表}
\setlength{\LTcapwidth}{\linewidth}
\begin{longtable}{@{}llp{0.42\textwidth}@{}}
\caption{常见二次曲面的标准方程与截痕特征}\\
\toprule
名称 & 标准方程 & 截痕特征 \\
\midrule
\endhead
\bottomrule
\endlastfoot
椭球面 & $\dfrac{x^{2}}{a^{2}}+\dfrac{y^{2}}{b^{2}}+\dfrac{z^{2}}{c^{2}}=1$ & 三组平行平面截得椭圆，有界（$|x|\le a$ 等） \\
单叶双曲面 & $\dfrac{x^{2}}{a^{2}}+\dfrac{y^{2}}{b^{2}}-\dfrac{z^{2}}{c^{2}}=1$ & 平行 $xOy$ 面截得椭圆且半轴随 $|h|$ 增大；平行 $yOz$ 面截得双曲线 \\
双叶双曲面 & $\dfrac{x^{2}}{a^{2}}-\dfrac{y^{2}}{b^{2}}-\dfrac{z^{2}}{c^{2}}=1$ & 分两叶（$|x|\ge a$）；$|x|>a$ 时截得椭圆 \\
椭圆抛物面 & $\dfrac{x^{2}}{2p}+\dfrac{y^{2}}{2q}=z\ (p,q\ \text{同号})$ & 平行 $xOy$ 面截得椭圆，平行 $yOz$ 面截得抛物线；无界 \\
双曲抛物面 & $\dfrac{x^{2}}{2p}-\dfrac{y^{2}}{2q}=z\ (p,q\ \text{同号})$ & 马鞍形；$z=h\neq0$ 截得双曲线，$z=0$ 截得两条相交直线 \\
二次锥面 & $\dfrac{x^{2}}{a^{2}}+\dfrac{y^{2}}{b^{2}}-\dfrac{z^{2}}{c^{2}}=0$ & 二次齐次；顶点在原点，$z=h\neq0$ 截得椭圆 \\
\end{longtable}
\endgroup

表中六个是真正“弯曲”的二次曲面；椭圆柱面、双曲柱面、抛物柱面这三类（也属二次曲面）的标准方程与母线方向见上一小节。椭球面、单叶双曲面、双曲抛物面、二次锥面还具有\textbf{直纹性}（曲面可由直线族生成），这是工程上（冷却塔、薄壳屋顶）常用的性质。用二次型的正交变换系统判别曲面类型，见第5章「拓展：空间中的曲面与二次曲面」一节。

\corepoint{
（1）旋转曲面公式中“与旋转轴同名的变量不动”，绕 $z$ 轴旋转时是 $y$ 换成 $\pm\sqrt{x^{2}+y^{2}}$，不是 $z$ 换；写错会得到完全不同的曲面。\\
（2）“方程只含两个变量 $\Rightarrow$ 柱面”这条判据里的“缺少的变量”决定母线方向：缺 $z$ 则母线平行 $z$ 轴。\\
（3）$\dfrac{x^{2}}{a^{2}}+\dfrac{y^{2}}{b^{2}}-\dfrac{z^{2}}{c^{2}}=1$ 与 $=0$ 只差一个常数项，图形却分别是单叶双曲面与二次锥面，判断时必须看清右端常数。\\
（4）截痕法只作定性判断；若要精确判定类型，应把二次型化为标准形，由特征值的符号与个数（惯性指数）分类。
}
